% Add to the manuscript preamble:
% \usepackage{graphicx}
% \usepackage{subcaption}
%
% From the repository/manuscript root:
% % Add to the manuscript preamble:
% \usepackage{graphicx}
% \usepackage{subcaption}
%
% From the repository/manuscript root:
% % Add to the manuscript preamble:
% \usepackage{graphicx}
% \usepackage{subcaption}
%
% From the repository/manuscript root:
% % Add to the manuscript preamble:
% \usepackage{graphicx}
% \usepackage{subcaption}
%
% From the repository/manuscript root:
% \input{figures/paper_heatmaps/maze_width070mm_seed0/heatmaps_individual.tex}
% If the assets are elsewhere, define \heatmapindividualdir before this input.
\providecommand{\heatmapindividualdir}{figures/paper_heatmaps/maze_width070mm_seed0/individual}

% The images contain no panel names or colorbars. LaTeX owns the subcaptions;
% the separately exported colorbar retains the same normalization as all maps.
% Keep all four panel widths equal. Do not crop/stretch individual colorbars.
% This 178 mm layout is intended to span both columns of a two-column paper.
% For a single-column manuscript use figure instead of figure*.

\begin{figure*}[t]
  \centering
  \captionsetup[subfigure]{font=small,labelformat=parens,labelsep=space,justification=centering}
  \includegraphics[width=0.60\textwidth]{\heatmapindividualdir/clearance_colorbar.pdf}\par
  \vspace{1mm}
  \begin{subfigure}[t]{0.21\textwidth}
    \centering
    \includegraphics[width=\linewidth]{\heatmapindividualdir/clearance_mpcc.pdf}
    \caption{MPCC}\label{fig:clearance-individual-mpcc}
  \end{subfigure}\hfill
  \begin{subfigure}[t]{0.21\textwidth}
    \centering
    \includegraphics[width=\linewidth]{\heatmapindividualdir/clearance_cc_ecs.pdf}
    \caption{CC-ECS}\label{fig:clearance-individual-cc-ecs}
  \end{subfigure}\hfill
  \begin{subfigure}[t]{0.21\textwidth}
    \centering
    \includegraphics[width=\linewidth]{\heatmapindividualdir/clearance_cc_sf.pdf}
    \caption{CC-SF}\label{fig:clearance-individual-cc-sf}
  \end{subfigure}\hfill
  \begin{subfigure}[t]{0.21\textwidth}
    \centering
    \includegraphics[width=\linewidth]{\heatmapindividualdir/clearance_cc_mf.pdf}
    \caption{CC-MF}\label{fig:clearance-individual-cc-mf}
  \end{subfigure}
  \caption{Clearance along the closed-loop trajectories in the maze for a
    differential-drive robot with a $150\times70$ mm rectangular footprint.
    Each method shows one trial with the same initial pose (seed 0).
    Clearance is the minimum of the exact footprint--obstacle distance and
    the signed map-boundary margin. All panels share the magma color scale
    over $0$--$40$ mm. A power-law color normalization ($\gamma=0.35$) is
    used to emphasize small clearances. Open circles and filled black stars
    mark the start and end positions, respectively.}
  \label{fig:clearance-individual}
\end{figure*}

\begin{figure*}[t]
  \centering
  \captionsetup[subfigure]{font=small,labelformat=parens,labelsep=space,justification=centering}
  \includegraphics[width=0.60\textwidth]{\heatmapindividualdir/velocity_colorbar.pdf}\par
  \vspace{1mm}
  \begin{subfigure}[t]{0.21\textwidth}
    \centering
    \includegraphics[width=\linewidth]{\heatmapindividualdir/velocity_mpcc.pdf}
    \caption{MPCC}\label{fig:velocity-individual-mpcc}
  \end{subfigure}\hfill
  \begin{subfigure}[t]{0.21\textwidth}
    \centering
    \includegraphics[width=\linewidth]{\heatmapindividualdir/velocity_cc_ecs.pdf}
    \caption{CC-ECS}\label{fig:velocity-individual-cc-ecs}
  \end{subfigure}\hfill
  \begin{subfigure}[t]{0.21\textwidth}
    \centering
    \includegraphics[width=\linewidth]{\heatmapindividualdir/velocity_cc_sf.pdf}
    \caption{CC-SF}\label{fig:velocity-individual-cc-sf}
  \end{subfigure}\hfill
  \begin{subfigure}[t]{0.21\textwidth}
    \centering
    \includegraphics[width=\linewidth]{\heatmapindividualdir/velocity_cc_mf.pdf}
    \caption{CC-MF}\label{fig:velocity-individual-cc-mf}
  \end{subfigure}
  \caption{Applied linear velocity along the closed-loop trajectories for
    the same maze trials (seed 0, differential-drive robot,
    $150\times70$ mm rectangular footprint).
    Each trajectory segment is colored by the linear velocity applied
    during that interval. All panels share a linear color scale over
    $0$--$0.18$ m/s. Open circles and filled black stars mark the start and
    end positions, respectively.}
  \label{fig:velocity-individual}
\end{figure*}

% Reference a whole figure: Fig.~\ref{fig:clearance-individual}.
% Reference a panel: Fig.~\ref{fig:clearance-individual}\subref{fig:clearance-individual-mpcc}.

% If the assets are elsewhere, define \heatmapindividualdir before this input.
\providecommand{\heatmapindividualdir}{figures/paper_heatmaps/maze_width070mm_seed0/individual}

% The images contain no panel names or colorbars. LaTeX owns the subcaptions;
% the separately exported colorbar retains the same normalization as all maps.
% Keep all four panel widths equal. Do not crop/stretch individual colorbars.
% This 178 mm layout is intended to span both columns of a two-column paper.
% For a single-column manuscript use figure instead of figure*.

\begin{figure*}[t]
  \centering
  \captionsetup[subfigure]{font=small,labelformat=parens,labelsep=space,justification=centering}
  \includegraphics[width=0.60\textwidth]{\heatmapindividualdir/clearance_colorbar.pdf}\par
  \vspace{1mm}
  \begin{subfigure}[t]{0.21\textwidth}
    \centering
    \includegraphics[width=\linewidth]{\heatmapindividualdir/clearance_mpcc.pdf}
    \caption{MPCC}\label{fig:clearance-individual-mpcc}
  \end{subfigure}\hfill
  \begin{subfigure}[t]{0.21\textwidth}
    \centering
    \includegraphics[width=\linewidth]{\heatmapindividualdir/clearance_cc_ecs.pdf}
    \caption{CC-ECS}\label{fig:clearance-individual-cc-ecs}
  \end{subfigure}\hfill
  \begin{subfigure}[t]{0.21\textwidth}
    \centering
    \includegraphics[width=\linewidth]{\heatmapindividualdir/clearance_cc_sf.pdf}
    \caption{CC-SF}\label{fig:clearance-individual-cc-sf}
  \end{subfigure}\hfill
  \begin{subfigure}[t]{0.21\textwidth}
    \centering
    \includegraphics[width=\linewidth]{\heatmapindividualdir/clearance_cc_mf.pdf}
    \caption{CC-MF}\label{fig:clearance-individual-cc-mf}
  \end{subfigure}
  \caption{Clearance along the closed-loop trajectories in the maze for a
    differential-drive robot with a $150\times70$ mm rectangular footprint.
    Each method shows one trial with the same initial pose (seed 0).
    Clearance is the minimum of the exact footprint--obstacle distance and
    the signed map-boundary margin. All panels share the magma color scale
    over $0$--$40$ mm. A power-law color normalization ($\gamma=0.35$) is
    used to emphasize small clearances. Open circles and filled black stars
    mark the start and end positions, respectively.}
  \label{fig:clearance-individual}
\end{figure*}

\begin{figure*}[t]
  \centering
  \captionsetup[subfigure]{font=small,labelformat=parens,labelsep=space,justification=centering}
  \includegraphics[width=0.60\textwidth]{\heatmapindividualdir/velocity_colorbar.pdf}\par
  \vspace{1mm}
  \begin{subfigure}[t]{0.21\textwidth}
    \centering
    \includegraphics[width=\linewidth]{\heatmapindividualdir/velocity_mpcc.pdf}
    \caption{MPCC}\label{fig:velocity-individual-mpcc}
  \end{subfigure}\hfill
  \begin{subfigure}[t]{0.21\textwidth}
    \centering
    \includegraphics[width=\linewidth]{\heatmapindividualdir/velocity_cc_ecs.pdf}
    \caption{CC-ECS}\label{fig:velocity-individual-cc-ecs}
  \end{subfigure}\hfill
  \begin{subfigure}[t]{0.21\textwidth}
    \centering
    \includegraphics[width=\linewidth]{\heatmapindividualdir/velocity_cc_sf.pdf}
    \caption{CC-SF}\label{fig:velocity-individual-cc-sf}
  \end{subfigure}\hfill
  \begin{subfigure}[t]{0.21\textwidth}
    \centering
    \includegraphics[width=\linewidth]{\heatmapindividualdir/velocity_cc_mf.pdf}
    \caption{CC-MF}\label{fig:velocity-individual-cc-mf}
  \end{subfigure}
  \caption{Applied linear velocity along the closed-loop trajectories for
    the same maze trials (seed 0, differential-drive robot,
    $150\times70$ mm rectangular footprint).
    Each trajectory segment is colored by the linear velocity applied
    during that interval. All panels share a linear color scale over
    $0$--$0.18$ m/s. Open circles and filled black stars mark the start and
    end positions, respectively.}
  \label{fig:velocity-individual}
\end{figure*}

% Reference a whole figure: Fig.~\ref{fig:clearance-individual}.
% Reference a panel: Fig.~\ref{fig:clearance-individual}\subref{fig:clearance-individual-mpcc}.

% If the assets are elsewhere, define \heatmapindividualdir before this input.
\providecommand{\heatmapindividualdir}{figures/paper_heatmaps/maze_width070mm_seed0/individual}

% The images contain no panel names or colorbars. LaTeX owns the subcaptions;
% the separately exported colorbar retains the same normalization as all maps.
% Keep all four panel widths equal. Do not crop/stretch individual colorbars.
% This 178 mm layout is intended to span both columns of a two-column paper.
% For a single-column manuscript use figure instead of figure*.

\begin{figure*}[t]
  \centering
  \captionsetup[subfigure]{font=small,labelformat=parens,labelsep=space,justification=centering}
  \includegraphics[width=0.60\textwidth]{\heatmapindividualdir/clearance_colorbar.pdf}\par
  \vspace{1mm}
  \begin{subfigure}[t]{0.21\textwidth}
    \centering
    \includegraphics[width=\linewidth]{\heatmapindividualdir/clearance_mpcc.pdf}
    \caption{MPCC}\label{fig:clearance-individual-mpcc}
  \end{subfigure}\hfill
  \begin{subfigure}[t]{0.21\textwidth}
    \centering
    \includegraphics[width=\linewidth]{\heatmapindividualdir/clearance_cc_ecs.pdf}
    \caption{CC-ECS}\label{fig:clearance-individual-cc-ecs}
  \end{subfigure}\hfill
  \begin{subfigure}[t]{0.21\textwidth}
    \centering
    \includegraphics[width=\linewidth]{\heatmapindividualdir/clearance_cc_sf.pdf}
    \caption{CC-SF}\label{fig:clearance-individual-cc-sf}
  \end{subfigure}\hfill
  \begin{subfigure}[t]{0.21\textwidth}
    \centering
    \includegraphics[width=\linewidth]{\heatmapindividualdir/clearance_cc_mf.pdf}
    \caption{CC-MF}\label{fig:clearance-individual-cc-mf}
  \end{subfigure}
  \caption{Clearance along the closed-loop trajectories in the maze for a
    differential-drive robot with a $150\times70$ mm rectangular footprint.
    Each method shows one trial with the same initial pose (seed 0).
    Clearance is the minimum of the exact footprint--obstacle distance and
    the signed map-boundary margin. All panels share the magma color scale
    over $0$--$40$ mm. A power-law color normalization ($\gamma=0.35$) is
    used to emphasize small clearances. Open circles and filled black stars
    mark the start and end positions, respectively.}
  \label{fig:clearance-individual}
\end{figure*}

\begin{figure*}[t]
  \centering
  \captionsetup[subfigure]{font=small,labelformat=parens,labelsep=space,justification=centering}
  \includegraphics[width=0.60\textwidth]{\heatmapindividualdir/velocity_colorbar.pdf}\par
  \vspace{1mm}
  \begin{subfigure}[t]{0.21\textwidth}
    \centering
    \includegraphics[width=\linewidth]{\heatmapindividualdir/velocity_mpcc.pdf}
    \caption{MPCC}\label{fig:velocity-individual-mpcc}
  \end{subfigure}\hfill
  \begin{subfigure}[t]{0.21\textwidth}
    \centering
    \includegraphics[width=\linewidth]{\heatmapindividualdir/velocity_cc_ecs.pdf}
    \caption{CC-ECS}\label{fig:velocity-individual-cc-ecs}
  \end{subfigure}\hfill
  \begin{subfigure}[t]{0.21\textwidth}
    \centering
    \includegraphics[width=\linewidth]{\heatmapindividualdir/velocity_cc_sf.pdf}
    \caption{CC-SF}\label{fig:velocity-individual-cc-sf}
  \end{subfigure}\hfill
  \begin{subfigure}[t]{0.21\textwidth}
    \centering
    \includegraphics[width=\linewidth]{\heatmapindividualdir/velocity_cc_mf.pdf}
    \caption{CC-MF}\label{fig:velocity-individual-cc-mf}
  \end{subfigure}
  \caption{Applied linear velocity along the closed-loop trajectories for
    the same maze trials (seed 0, differential-drive robot,
    $150\times70$ mm rectangular footprint).
    Each trajectory segment is colored by the linear velocity applied
    during that interval. All panels share a linear color scale over
    $0$--$0.18$ m/s. Open circles and filled black stars mark the start and
    end positions, respectively.}
  \label{fig:velocity-individual}
\end{figure*}

% Reference a whole figure: Fig.~\ref{fig:clearance-individual}.
% Reference a panel: Fig.~\ref{fig:clearance-individual}\subref{fig:clearance-individual-mpcc}.

% If the assets are elsewhere, define \heatmapindividualdir before this input.
\providecommand{\heatmapindividualdir}{figures/paper_heatmaps/maze_width070mm_seed0/individual}

% The images contain no panel names or colorbars. LaTeX owns the subcaptions;
% the separately exported colorbar retains the same normalization as all maps.
% Keep all four panel widths equal. Do not crop/stretch individual colorbars.
% This 178 mm layout is intended to span both columns of a two-column paper.
% For a single-column manuscript use figure instead of figure*.

\begin{figure*}[t]
  \centering
  \captionsetup[subfigure]{font=small,labelformat=parens,labelsep=space,justification=centering}
  \includegraphics[width=0.60\textwidth]{\heatmapindividualdir/clearance_colorbar.pdf}\par
  \vspace{1mm}
  \begin{subfigure}[t]{0.21\textwidth}
    \centering
    \includegraphics[width=\linewidth]{\heatmapindividualdir/clearance_mpcc.pdf}
    \caption{MPCC}\label{fig:clearance-individual-mpcc}
  \end{subfigure}\hfill
  \begin{subfigure}[t]{0.21\textwidth}
    \centering
    \includegraphics[width=\linewidth]{\heatmapindividualdir/clearance_cc_ecs.pdf}
    \caption{CC-ECS}\label{fig:clearance-individual-cc-ecs}
  \end{subfigure}\hfill
  \begin{subfigure}[t]{0.21\textwidth}
    \centering
    \includegraphics[width=\linewidth]{\heatmapindividualdir/clearance_cc_sf.pdf}
    \caption{CC-SF}\label{fig:clearance-individual-cc-sf}
  \end{subfigure}\hfill
  \begin{subfigure}[t]{0.21\textwidth}
    \centering
    \includegraphics[width=\linewidth]{\heatmapindividualdir/clearance_cc_mf.pdf}
    \caption{CC-MF}\label{fig:clearance-individual-cc-mf}
  \end{subfigure}
  \caption{Clearance along the closed-loop trajectories in the maze for a
    differential-drive robot with a $150\times70$ mm rectangular footprint.
    Each method shows one trial with the same initial pose (seed 0).
    Clearance is the minimum of the exact footprint--obstacle distance and
    the signed map-boundary margin. All panels share the magma color scale
    over $0$--$40$ mm. A power-law color normalization ($\gamma=0.35$) is
    used to emphasize small clearances. Open circles and filled black stars
    mark the start and end positions, respectively.}
  \label{fig:clearance-individual}
\end{figure*}

\begin{figure*}[t]
  \centering
  \captionsetup[subfigure]{font=small,labelformat=parens,labelsep=space,justification=centering}
  \includegraphics[width=0.60\textwidth]{\heatmapindividualdir/velocity_colorbar.pdf}\par
  \vspace{1mm}
  \begin{subfigure}[t]{0.21\textwidth}
    \centering
    \includegraphics[width=\linewidth]{\heatmapindividualdir/velocity_mpcc.pdf}
    \caption{MPCC}\label{fig:velocity-individual-mpcc}
  \end{subfigure}\hfill
  \begin{subfigure}[t]{0.21\textwidth}
    \centering
    \includegraphics[width=\linewidth]{\heatmapindividualdir/velocity_cc_ecs.pdf}
    \caption{CC-ECS}\label{fig:velocity-individual-cc-ecs}
  \end{subfigure}\hfill
  \begin{subfigure}[t]{0.21\textwidth}
    \centering
    \includegraphics[width=\linewidth]{\heatmapindividualdir/velocity_cc_sf.pdf}
    \caption{CC-SF}\label{fig:velocity-individual-cc-sf}
  \end{subfigure}\hfill
  \begin{subfigure}[t]{0.21\textwidth}
    \centering
    \includegraphics[width=\linewidth]{\heatmapindividualdir/velocity_cc_mf.pdf}
    \caption{CC-MF}\label{fig:velocity-individual-cc-mf}
  \end{subfigure}
  \caption{Applied linear velocity along the closed-loop trajectories for
    the same maze trials (seed 0, differential-drive robot,
    $150\times70$ mm rectangular footprint).
    Each trajectory segment is colored by the linear velocity applied
    during that interval. All panels share a linear color scale over
    $0$--$0.18$ m/s. Open circles and filled black stars mark the start and
    end positions, respectively.}
  \label{fig:velocity-individual}
\end{figure*}

% Reference a whole figure: Fig.~\ref{fig:clearance-individual}.
% Reference a panel: Fig.~\ref{fig:clearance-individual}\subref{fig:clearance-individual-mpcc}.
